\documentclass[11pt]{article}

\usepackage[margin=1in]{geometry}
\usepackage{amsmath,amssymb,amsthm}
\usepackage{booktabs}
\usepackage{graphicx}
\usepackage{microtype}
\usepackage{url}
\usepackage[breaklinks,colorlinks=true,linkcolor=blue,citecolor=blue,urlcolor=blue]{hyperref}

\newtheorem{definition}{Definition}
\newtheorem{proposition}{Proposition}
\newcommand{\code}[1]{\nolinkurl{#1}}
\newcommand{\WM}{\mathrm{WM}}
\newcommand{\FCA}{\mathrm{FCA}}
\newcommand{\Pow}{\mathcal{P}}
\newcommand{\lowerrel}{\underline{I}}
\newcommand{\upperrel}{\overline{I}}
\newcommand{\yes}{\mathsf{T}}
\newcommand{\no}{\mathsf{F}}
\newcommand{\unknown}{\mathsf{U}}
\newcommand{\missing}{\mathsf{M}}

% Generated by generate_wm_fca_benchmark_tables.py; do not edit.
\newcommand{\WMFCACaseCount}{405}
\newcommand{\WMFCAAssertionCount}{12,150}
\newcommand{\WMFCAEventPrefixCount}{30}
\newcommand{\WMFCAExactContextCount}{28}
\newcommand{\WMFCAExactConceptCount}{3,275}
\newcommand{\WMFCAExactBasisCount}{1,330}
\newcommand{\WMFCAExactCoverCount}{11,734}
\newcommand{\WMFCAQueryContextCount}{19}
\newcommand{\WMFCAQueryCount}{1,000}
\newcommand{\WMFCACorruptionContextCount}{9}
\newcommand{\WMFCASeedCount}{5}
\newcommand{\WMFCAMissingRates}{0.10, 0.30, 0.50}
\newcommand{\WMFCAFlipRates}{0.00, 0.05, 0.10}
\newcommand{\WMFCAUnknownRate}{0.05}
\newcommand{\WMFCAProtocolHash}{\texttt{7dcd2936b014...}}
\newcommand{\WMFCAReceiptHash}{\texttt{835ac74af9be...}}
\newcommand{\WMFCAReceiptContentHash}{\texttt{40c755197680...}}
\newcommand{\WMFCAIncompleteProtocolHash}{\texttt{cc0596bbeb8e...}}
\newcommand{\WMFCAIncompleteReceiptHash}{\texttt{a414d827b9d3...}}
\newcommand{\WMFCAIncompleteReceiptContentHash}{\texttt{003a0e25831c...}}
\newcommand{\WMFCAErgonomicsReceiptHash}{\texttt{6fdc93a90a89...}}
\newcommand{\WMFCAErgonomicsPeakMiB}{19.5}
\newcommand{\WMFCALibraryLines}{2,309}
\newcommand{\WMFCALibraryBytes}{79,938}
\newcommand{\WMFCALibraryTypeDeclarations}{110}
\newcommand{\WMFCALibraryEquations}{249}
\newcommand{\WMFCACertificateCoverageCount}{14}
\newcommand{\WMFCACalibrationStratumCount}{6}

\newcommand{\WMFCAExactValidationRows}{%
Published FCA examples & 6 & 8--35 $\times$ 7--14 & 344 & 104 & 887 \\
Mizar family & 13 & 27--143 $\times$ 1--32 & 767 & 986 & 1,729 \\
FCApy controlled random & 9 & 10--100 $\times$ 10 & 2,164 & 240 & 9,118 \\
}

\newcommand{\WMFCAQueryValidationRows}{%
FCApy controlled random & 18 & 100 $\times$ 50 & 918 \\
Bob Ross & 1 & 403 $\times$ 67 & 82 \\
}

\newcommand{\WMFCAErgonomicsRows}{%
Bob Ross query batch & 82 closure queries & 9.989 & 8,209 & yes \\
Immutable event replay & 14 event operations & 12.868 & 1,088 & yes \\
Certificate checks & 24 verified assertions & 13.266 & 1,809 & yes \\
}

\newcommand{\WMFCAModelResultRows}{%
Crisp single source & 0.723 & 0.742 & 0.153 & 0.138 & 0.368 \\
Raw majority (2) & 0.723 & 0.742 & 0.153 & 0.138 & 0.368 \\
Source majority (2) & 0.723 & 0.742 & 0.153 & 0.138 & 0.368 \\
WM permissive (1 group) & 0.838 & 0.772 & 0.266 & 0.215 & 0.518 \\
WM robust (2 groups) & 0.569 & 0.704 & 0.139 & 0.126 & 0.250 \\
WM exact consensus & 0.830 & 0.766 & 0.222 & 0.204 & 0.485 \\
}

\newcommand{\WMFCAIncompleteBaselineRows}{%
Incomplete lower completion & 0.723 & 0.742 & 0.153 & 0.138 & 0.368 \\
Incomplete upper completion & 0.736 & 0.675 & 0.146 & 0.050 & 0.337 \\
WM permissive (1 group) & 0.838 & 0.772 & 0.266 & 0.215 & 0.518 \\
}

\newcommand{\WMFCAPermissiveDeltaRows}{%
Cell F1 & +0.114 & 370/1/34 \\
Closure exact & +0.030 & 168/158/79 \\
Canonical-basis F1 & +0.113 & 212/109/84 \\
Concept F1 & +0.076 & 216/41/148 \\
Cover-edge F1 & +0.150 & 324/18/63 \\
}

\newcommand{\WMFCAEnvelopeRows}{%
Cells & 0.799 & 0.919 & +0.120 & 0.239 \\
Closures & 0.742 & 0.845 & +0.103 & 0.118 \\
}

\newcommand{\WMFCACalibrationRows}{%
Crisp hard point & $\{0,1\}$ & 0.201 & 0.201 & 0.799 & 0.000 \\
Incomplete completion midpoint & $\{0,0.5,1\}$ & 0.116 & 0.126 & 0.968 & 0.334 \\
WM model-family midpoint & $\{0,0.5,1\}$ & 0.141 & 0.179 & 0.919 & 0.239 \\
}

\newcommand{\WMFCADuplicateInvariantRate}{1.000}
\newcommand{\WMFCAGroupCopyInvariantRate}{1.000}
\newcommand{\WMFCARawCopyInvariantRate}{0.007}


\title{World-Model Formal Concept Analysis:\\
Exact Recovery, Provenance-Safe Revision, and Contextual Envelopes}
\author{Zar Goertzel}
\date{July 2026}

\begin{document}
\maketitle

\begin{abstract}
Formal Concept Analysis (FCA) extracts concepts, implications, and a concept
lattice from a binary object--attribute relation. Real knowledge stores are
usually less obliging: an incidence can be false, unobserved, explicitly
indeterminate, duplicated by a pipeline, revised, or later forgotten. We
present a world-model (WM) formulation in which the primary object is a finite
state of provenance-stamped observations and classical FCA is an exact late
projection. The implementation runs natively in CeTTa and supports immutable
revision, forgetting, cell-level explanations, source/dependence-group gates,
bounded concept and implication enumeration, and robust/permissive incidence
envelopes. A Lean~4 development proves exact order-isomorphic recovery of
classical FCA, duplicate-idempotent revision, disjoint revision/forgetting
laws, dependence-group invariance, and contextual-envelope obligations without
\code{sorry}, wishlist stubs, or new axioms.

We first establish exact parity on \WMFCAExactContextCount{} finite contexts,
covering published examples, a 13-context Mizar family, and FCApy controlled
random contexts. The checked artifacts comprise \WMFCAExactConceptCount{}
concepts, \WMFCAExactBasisCount{} canonical implications, and
\WMFCAExactCoverCount{} cover edges. A further \WMFCAQueryContextCount{} wider
contexts, including the $403\times67$ Bob Ross matrix, supply
\WMFCAQueryCount{} exact extent/closure queries without unbounded lattice
enumeration. We then run a frozen \WMFCACaseCount{}-case corruption study.
Dependence-aware WM projections are invariant to correlated copies
in every case, while raw/source counting is invariant in only a
\WMFCARawCopyInvariantRate{} fraction of cases. A permissive WM point
projection improves the descriptive mean recovery of cells, concepts,
implications, and covers over the crisp single-source projection; a robust
projection does worse, showing that the WM degrees of freedom are substantive
rather than uniformly beneficial.
The robust/permissive envelope raises cell and closure coverage at the explicit
cost of nonzero interval width; its midpoint also has lower Brier score and ECE
than the crisp hard forecast in every frozen corruption stratum. This scores
the midpoint as a point forecast; it does not establish calibration of the
set-valued interval.
In a separately frozen incomplete-context comparison, the permissive WM
projection achieves better descriptive-mean point and structural recovery than
both the lower and upper completions. The completion interval, however, has higher
cell coverage and lower midpoint Brier score and ECE than the narrower WM
envelope. Thus the experiment exposes a coverage--width tradeoff rather than a
uniform winner.
The study is descriptive, not an inferential significance claim, and uses
synthetic corruption rather than a natural-domain uncertainty corpus.
\end{abstract}

\section{Introduction}

Formal Concept Analysis begins with a formal context $K=(G,M,I)$: objects $G$,
attributes $M$, and a binary incidence relation $I\subseteq G\times M$
\cite{wille1982,ganterwille1999}. From the Galois connection
\[
  A'=\{m\in M:\forall g\in A,\ gIm\},\qquad
  B'=\{g\in G:\forall m\in B,\ gIm\},
\]
it derives concepts $(A,B)$ satisfying $A'=B$ and $B'=A$, attribute
implications, and a complete concept lattice. This is an exact and useful
theory once $I$ is fixed. The difficult step in an evidential system is often
the condition \emph{once $I$ is fixed}.

A missing database cell is not a negative observation. Two reports can be
nominally distinct but descend from the same experiment. An observation can be
retracted without invalidating unrelated evidence. Different admissible
extraction policies can induce different crisp relations, and a query may need
the range across those policies rather than one prematurely selected matrix.
Encoding all of these cases directly as Boolean incidence loses information
before concept formation begins.

We therefore put a small world-model calculus in front of classical FCA.
States, revision, forgetting, query, and explanation are primary. A gate
projects one evidence profile into a crisp relation only when an exact FCA
artifact is requested. A family of admissible models yields lower (robust) and
upper (permissive) relations. The formulation is deliberately not restricted
to PLN truth values: counts, four-valued observations, interval evidence, or
other ordered evidence carriers can instantiate the interface.

The contributions are:
\begin{enumerate}
  \item a finite provenance-state semantics with immutable revision and
        forgetting, distinct false, unknown, and missing statuses, and
        dependence-aware extraction;
  \item native CeTTa operations for exact FCA projection, bounded
        NextClosure-style enumeration, canonical implications, cover edges,
        packed scale contexts, and cell-level provenance explanations;
  \item Lean proofs connecting the state semantics to classical FCA and
        establishing revision, forgetting, deduplication, dependence, and
        contextual-envelope laws;
  \item exact parity against independent FCA oracles on
        \WMFCAExactContextCount{} enumeration contexts and
        \WMFCAQueryContextCount{} wider query contexts; and
  \item a frozen corruption study that separates point-recovery quality,
        interval coverage/width, and invariance to duplicated evidence,
        together with a separately frozen lower/upper-completion FCA
        comparison.
\end{enumerate}

The central result is not that one WM gate is always superior. It is that
classical FCA is recovered exactly when the state is complete and the model
family is a singleton, while richer state and model structure can be retained
and tested when those assumptions do not hold.

\section{From observations to formal contexts}
\label{sec:semantics}

\subsection{Provenance states}

Let $S$ be a finite type of stamps. A provenance catalog fixes, for every
$s\in S$, its source, declared dependence group, and query-local evidence
contribution:
\[
  C_{\mathrm{src}}:S\to\mathit{Source},\qquad
  C_{\mathrm{grp}}:S\to\mathit{Group},\qquad
  C_{\mathrm{obs}}:S\times G\times M\to E .
\]
A stamp denotes immutable content. Reusing a stamp with conflicting source,
group, or observation is therefore unrepresentable in the formal model and
rejected by the runtime.

A state $W\in\Pow_{\mathrm{fin}}(S)$ is the set of active stamps. For an
additive evidence carrier $E$,
\[
  e_W(g,m)=\sum_{s\in W} C_{\mathrm{obs}}(s,g,m).
\]
Revision and forgetting are
\[
  \operatorname{Revise}(W,\Delta)=W\cup\Delta,\qquad
  \operatorname{Forget}(D,W)=W\setminus D.
\]
Union makes exact replay idempotent. Additivity is asserted only for disjoint
stamp sets:
\[
  W\cap\Delta=\varnothing
  \Longrightarrow
  e_{W\cup\Delta}=e_W+e_\Delta .
\]
This guard is important. Treating idempotent union as an additive monoid would
silently count a replayed observation twice.

The executable carrier specializes observations to four statuses:
$\{\yes,\no,\unknown,\missing\}$. $\no$ is an observed negative;
$\unknown$ is an explicit indeterminate report; $\missing$ means the state
contains no applicable observation. Explanations return the contributing
stamps with source and group metadata rather than only the projected bit.

\subsection{Gates and exact FCA projection}

An evidence gate $q:E\to\{\mathrm{false},\mathrm{true}\}$ induces
\[
  I_{W,q}(g,m)\quad\Longleftrightarrow\quad q(e_W(g,m)).
\]
The implementation includes exact four-status extraction and thresholds over
positive observations, distinct sources, and distinct dependence groups.
Classical FCA is then applied to $I_{W,q}$, not redefined. In particular, the
Lean map \code{formedConceptOrderIsoCrispConcept} is an order isomorphism
between exact formed concepts and the standard FCA concept carrier for the
same gated relation. Thus the WM lane is a conservative extension, not a
competing notion of crisp concept.

Dependence groups alter what is counted. Let
\[
  \Gamma_W(g,m)=
  \{C_{\mathrm{grp}}(s):s\in W,\ C_{\mathrm{obs}}(s,g,m)\neq0\}.
\]
The group-count statistic is $|\Gamma_W(g,m)|$. If every group contributed by
$\Delta$ is already represented in $W$ at every cell, then
\[
  |\Gamma_{W\cup\Delta}(g,m)|=|\Gamma_W(g,m)|
\]
and every FCA family extracted through a group-count threshold is unchanged.
This is a semantic guarantee about a declared dependence partition; it does
not prove that the declaration itself is scientifically correct.

\subsection{Contextual envelopes}

A contextual model $\theta=(\mu_\theta,q_\theta)$ contains an evidence profile
and a gate. For a nonempty admissible family $\Theta$, define
\[
\begin{aligned}
  \lowerrel(g,m)
    &\Longleftrightarrow \forall\theta\in\Theta,\
      q_\theta(\mu_\theta(g,m)),\\
  \upperrel(g,m)
    &\Longleftrightarrow \exists\theta\in\Theta,\
      q_\theta(\mu_\theta(g,m)).
\end{aligned}
\]
Then $\lowerrel\subseteq\upperrel$. We call the pair an incidence envelope.
It supports robust and permissive FCA queries without treating an unobserved
cell as false. For $\Theta=\{\theta\}$, both relations reduce exactly to
$I_\theta$, and both concept-family projections reduce to ordinary FCA on that
relation.

The envelope is intentionally a relation-level construction. The concept
family of $\lowerrel$ is not generally the set-theoretic intersection of the
concept families of all $\theta$, and similarly for $\upperrel$. Making the
aggregation point explicit prevents an invalid interchange between Galois
closure and model-family quantification.

\section{Implementation and formal assurance}
\label{sec:implementation}

\subsection{Native CeTTa execution}

The public implementation is a CeTTa library consumed by a benchmark suite.
Small contexts use transparent MeTTa definitions as differential
specifications. Scale contexts use validated packed binary rows and native
operations for:
\begin{itemize}
  \item row and event-state construction;
  \item status, evidence, source, group, and explanation queries;
  \item immutable revise/remember/forget operations;
  \item column-index construction and batched extent/closure queries;
  \item replayable closure, implication, and concept certificates, retaining
        the indexed state, gate, query, intermediate extent, and closure, and
        exposing the exact finite boundary of consulted cells; and
  \item closed-intent enumeration, Duquenne--Guigues canonical bases, and
        Hasse cover edges under explicit output/step bounds.
\end{itemize}

The bounds are part of the semantic contract: exhausting a bound fails loudly
rather than returning a prefix as a complete lattice. Contexts with at most
10 attributes are assigned to exact enumeration in the scale protocol. Wider
contexts are query-only. This is not a claim that a general MeTTa runtime
should outperform specialized FCA software.

An earlier recursive state/query path was replaced by bounded native carriers.
No benchmark in this paper uses the retired expansion of a large packed matrix
into one interpreted term per cell. Malformed rows, invalid statuses,
out-of-universe observations, conflicting stamp reuse, and exceeded bounds
have negative tests in addition to ordinary positive examples.

\subsection{Lean obligations}

The new Lean module builds on the existing FCA recovery bridge. Its primary
definitions are a finite stamp set, immutable catalog, union revision, set
difference forgetting, group support, exact concept-family extraction, and a
contextual envelope. The audited theorems include:
\begin{itemize}
  \item \code{revise_self}, \code{revise_comm}, and \code{revise_assoc};
  \item \code{forget_revision_of_disjoint} and
        \code{evidence_forget_revision_of_disjoint};
  \item \code{evidence_revise_of_disjoint}, exposing rather than hiding the
        freshness condition;
  \item \code{groupCountConceptFamily_revise_eq_left_of_groups_subset};
  \item \code{finiteConceptFamily_forget_revision_of_disjoint};
  \item the robust-to-permissive implication, singleton reductions, and
        contextual-envelope validity; and
  \item positive and negative executable canaries distinguishing a correlated
        stamp, an independent group, and an exact duplicate.
\end{itemize}

The module contains no \code{sorry}, \code{theorem_wanted}, or declared axiom.
Kernel dependency inspection of the load-bearing theorems reports only Lean's
standard \code{propext}, \code{Classical.choice}, and \code{Quot.sound} where
needed. A native-evaluation canary was deliberately replaced by an ordinary
kernel proof, so no native-decision trust axiom remains in the audited claims.
These theorems formalize the abstract state and projection laws; they are not a
formal verification of the C runtime.

\section{Exact validation}
\label{sec:exact}

\subsection{Datasets and independent oracles}

The exact suite contains three families:
\begin{enumerate}
  \item six published FCA examples: Planets, Living Beings, Animals, Binary
        Relations, Tea Lady, and Music;
  \item 13 contexts extracted from a pinned Mizar family; and
  \item nine 10-attribute controlled-random contexts reproduced from the
        FCApy benchmark generator \cite{fcapy}.
\end{enumerate}
Planets is checked against the published \texttt{fcaR} concepts and
implications \cite{fcar}. The remaining exact contexts are reconstructed by a
lectic NextClosure implementation, an independently ordered level-wise
implementation, and the \texttt{concepts} FCA package; lattice neighbors are
also compared. Inputs and dependency versions are content-addressed.

\begin{table}[t]
\centering
\small
\caption{Exact-enumeration validation. Artifact columns are totals within a
family; every concept, canonical implication, and cover edge is compared as a
set, not only by count.}
\label{tab:exact}
\begin{tabular}{lrrrrr}
\toprule
Family & Contexts & Object $\times$ attribute range &
Concepts & Basis & Covers\\
\midrule
\WMFCAExactValidationRows
\bottomrule
\end{tabular}
\end{table}

The \WMFCAQueryContextCount{} wider contexts, including the FiveThirtyEight Bob Ross matrix
\cite{bobross}, are not silently omitted. They have a separately frozen
query protocol containing the empty and full attribute sets, every singleton,
fixed prefixes, and fixed adjacent pairs. Direct set-intersection oracles and
\texttt{concepts} agree on all extents and closures.

\begin{table}[t]
\centering
\small
\caption{Query-only scale validation. Maximum size is the largest context in
the family; no complete lattice is claimed for these cases.}
\label{tab:query}
\begin{tabular}{lrrr}
\toprule
Family & Contexts & Maximum size & Exact queries\\
\midrule
\WMFCAQueryValidationRows
\bottomrule
\end{tabular}
\end{table}

\subsection{Revision and forgetting}

The event protocol contains 14 operations: forgetting and restoring a packed
stamp, repeating the same stamp, adding and removing a conflict, adding and
removing explicit unknown evidence, corroboration, a multi-observation
revision, and scoped forgetting by group, source, context, and cell. Both a
$30\times30$ random context and Bob Ross are checked at the initial state and
after every event, yielding \WMFCAEventPrefixCount{} prefix states. At each
prefix,
\[
  Q(\operatorname{foldUpdates}(W_0,e_1,\ldots,e_t))
  =
  Q(\operatorname{buildFresh}(\operatorname{survivors}_t))
\]
for the frozen status, evidence, provenance, extent, and closure queries.
Independent direct queries also agree with \texttt{concepts}. This validates
incremental semantics against fresh recomputation; it does not prove an
asymptotic update advantage.

\subsection{Staged ergonomics measurement}

A small receipt-backed campaign measures only fixtures whose golden outputs
were already validated. Each case has one warm-up and five measured process invocations; an
invocation is retained only when its output is byte-identical to the expected
output. Timings include process startup, parsing, native execution, and golden
validation. They are therefore end-to-end observations on one machine, not
kernel-only timings or comparisons with specialized FCA engines.

\begin{table}[ht]
\centering
\small
\caption{Descriptive native-runtime ergonomics. Rates have different work
units by row and should not be compared across rows.}
\label{tab:ergonomics}
\begin{tabular}{llrrc}
\toprule
Case & Work per invocation & Median ms & Units/s & Golden\\
\midrule
\WMFCAErgonomicsRows
\bottomrule
\end{tabular}
\end{table}

The recorded campaign peak child RSS is \WMFCAErgonomicsPeakMiB{} MiB. The
public WM/FCA library has \WMFCALibraryLines{} lines
(\WMFCALibraryBytes{} bytes), \WMFCALibraryTypeDeclarations{} type
declarations, and \WMFCALibraryEquations{} equations. The certificate gate
covers \WMFCACertificateCoverageCount{} named positive, negative, tamper, cell
boundary, history, and typed-wrapper behaviors. These counts measure source
and test surface, not semantic complexity. The complete environment, input,
binary, runner, protocol, program, expected-output, and library hashes are in
receipt \WMFCAErgonomicsReceiptHash.

\section{Frozen corruption study}
\label{sec:corruption}

\subsection{Protocol}

The protocol was frozen before the complete run and is identified by
\WMFCAProtocolHash. It crosses:
\begin{itemize}
  \item \WMFCACorruptionContextCount{} FCApy contexts with 10 attributes and
        10, 30, or 100 objects at densities 0.1, 0.5, and 0.9;
  \item \WMFCASeedCount{} deterministic seeds;
  \item missingness rates $\{\WMFCAMissingRates\}$;
  \item bit-flip rates $\{\WMFCAFlipRates\}$; and
  \item an explicit unknown rate of \WMFCAUnknownRate.
\end{itemize}
The Cartesian product contains \WMFCACaseCount{} cases. Each case has a
primary layer, an exact matrix copy from a distinct nominal source in the same
dependence group, and an independently corrupted layer in another group.

Six models were fixed:
\begin{description}
  \item[Crisp single source] at least one positive in the primary layer.
  \item[Raw majority (2)] at least two positive observation packets.
  \item[Source majority (2)] at least two nominal positive sources.
  \item[WM permissive] at least one positive dependence group.
  \item[WM robust] at least two positive dependence groups.
  \item[WM exact consensus] at least one positive and no observed negative
        across the selected layers.
\end{description}
No model was selected or tuned after inspecting outcomes. For every case,
CeTTa evaluates the layer projections and checks them against the independent
orchestrator, giving \WMFCAAssertionCount{} native semantic assertions.

Point projections are scored against the uncorrupted context by cell F1,
held-out closure exactness, concept-set F1, canonical-basis F1, and cover-edge
F1. The contextual family is scored by cell and closure coverage together
with interval width. Exact duplicate and correlated-copy invariance are
separate metrics.

\paragraph{Incomplete-context extension.}
After the main WM study, but before executing this comparator, we froze a
parameter-free lower/upper-completion extension over the same
\WMFCACaseCount{} source cases. In the epistemic incomplete-context view, each
cell is positive, negative, or unknown, and a completion replaces unknown
values by definite ones \cite{chacongomez2026}. We use only the primary
evidence layer: the lower completion maps unknown and missing cells to false,
and the upper completion maps them to true; both preserve known positive and
negative observations. The source layer is reconstructed and must match the
original receipt exactly before it is scored. Observed bit flips remain
observations in both completions.

Both completion relations are scored separately on the same point and
structural metrics. We do not treat their concept families, implication bases,
or closures as set-valued intervals: relation inclusion does not in general
make those derived families nested. The extension has no fitted parameter,
model selection, or post-result tuning and is identified by
\WMFCAIncompleteProtocolHash.

\subsection{Point recovery}

\begin{table}[t]
\centering
\scriptsize
\caption{Descriptive means across all \WMFCACaseCount{} frozen cases. The
robust model is substantially worse; the permissive and exact-consensus models
are better on these averages. No significance test is implied.}
\label{tab:model-results}
\resizebox{\textwidth}{!}{%
\begin{tabular}{lrrrrr}
\toprule
Model & Cell F1 & Closure exact & Basis F1 & Concept F1 & Cover F1\\
\midrule
\WMFCAModelResultRows
\bottomrule
\end{tabular}}
\end{table}

The raw- and source-majority rows equal the crisp row because the duplicated
primary matrix supplies their second positive vote. This equality is not a
desirable robustness result: removing the correlated copy changes these
projections in almost every case.

\begin{table}[t]
\centering
\scriptsize
\caption{Frozen incomplete-context comparison on the same
\WMFCACaseCount{} cases. Each completion is evaluated as its own crisp
relation; no concept-family interval is implied.}
\label{tab:incomplete-baselines}
\resizebox{\textwidth}{!}{%
\begin{tabular}{lrrrrr}
\toprule
Model & Cell F1 & Closure exact & Basis F1 & Concept F1 & Cover F1\\
\midrule
\WMFCAIncompleteBaselineRows
\bottomrule
\end{tabular}}
\end{table}

The lower completion is exactly the crisp single-source projection: both
retain observed positives and map negative, unknown, and missing cells to
non-incidence. A standard incomplete-context construction thus independently
recovers the conservative baseline. The WM permissive projection has
higher descriptive mean recovery than both completions on every metric in
Table~\ref{tab:incomplete-baselines}; the next subsection shows that this point
advantage does not dominate the interval comparison.

Table~\ref{tab:permissive-delta} shows paired descriptive differences for the
permissive group model. A win/tie/loss is computed per frozen case, so the
positive means do not hide that concept and implication recovery sometimes
worsens.

\begin{table}[t]
\centering
\small
\caption{WM permissive minus crisp single-source. W/T/L denotes paired cases.}
\label{tab:permissive-delta}
\begin{tabular}{lrr}
\toprule
Metric & Mean difference & W/T/L\\
\midrule
\WMFCAPermissiveDeltaRows
\bottomrule
\end{tabular}
\end{table}

\subsection{Envelopes and provenance invariance}

The three WM incidence projections (permissive, robust, and exact consensus)
form the evaluated envelope. Coverage is the fraction of true cell labels or
true closure memberships contained in the lower/upper interval. Width is the
corresponding average binary width (normalized for closures).

\begin{table}[t]
\centering
\small
\caption{Contextual-envelope coverage and sharpness. Higher coverage is
purchased with explicit width; this table is not a point-accuracy comparison.}
\label{tab:envelope}
\begin{tabular}{lrrrr}
\toprule
Target & Crisp coverage & WM coverage & Gain & Mean width\\
\midrule
\WMFCAEnvelopeRows
\bottomrule
\end{tabular}
\end{table}

The coverage gain is positive in every frozen missingness and flip-rate
stratum, but a wider interval can achieve coverage mechanically. The
scientific object is therefore the coverage--width pair, not coverage alone.
This experiment does not establish optimal sharpness or calibration of the
set-valued interval.

The main protocol also preregistered Brier score and expected calibration error
(ECE). For this point-scoring view only, an incidence interval $[l,u]$ is
mapped to its midpoint $(l+u)/2\in\{0,0.5,1\}$. The separately frozen
completion comparator uses the same map, while the crisp baseline emits hard
zero/one forecasts. Table~\ref{tab:calibration} averages the case-level scores;
lower Brier and ECE, higher coverage, and lower width are desirable but need
not be attained by the same method.

\begin{table}[ht]
\centering
\small
\caption{Cell-level point forecasts and interval coverage/width. Brier and ECE
score interval midpoints; they do not establish calibration of the set-valued
intervals.}
\label{tab:calibration}
\begin{tabular}{lrrrrr}
\toprule
Forecast & Support & Brier & ECE & Coverage & Width\\
\midrule
\WMFCACalibrationRows
\bottomrule
\end{tabular}
\end{table}

The WM-minus-crisp differences in midpoint Brier score and ECE are
negative---the midpoint forecast scores better---in each of the
\WMFCACalibrationStratumCount{} separately frozen missingness and flip-rate
strata. This is descriptive evidence that the disagreement state carries
useful uncertainty information. It is not an inferential comparison, and
midpoint calibration does not imply coverage calibration for the interval.
The lower/upper-completion midpoint has lower Brier score and ECE and its
interval covers more cells than the WM family, but it is also wider. The fixed
comparison therefore supports no claim that either interval uniformly
dominates the other.

Every model is invariant to replaying an evidence layer with the same identity:
the observed rate is \WMFCADuplicateInvariantRate. The two dependence-group
models are invariant to removing the distinct-source correlated copy at rate
\WMFCAGroupCopyInvariantRate. Raw and source-counting thresholds have rate
\WMFCARawCopyInvariantRate. This ablation isolates the benefit of preserving a
declared dependence structure: nominal source multiplicity is not allowed to
masquerade as independent support.

\section{What the WM layer contributes}
\label{sec:discussion}

\paragraph{Conservative exactness.}
When observations are complete and one extraction model is fixed, WM/FCA
reconstructs ordinary FCA exactly. This is supported both by the Lean order
isomorphism and by full artifact parity over the exact suite.

\paragraph{Reversible evidence state.}
Every cell judgment can be traced to active stamps, and forgetting has a
precise meaning: recompute from the surviving evidence. The lifecycle tests
exercise negative, unknown, missing, conflicting, duplicated, and
foreign-context observations rather than only monotone addition. Query-specific
closure certificates retain both stages of the Galois derivation: the queried
attributes, their extent, and the resulting closure. Implication and concept
certificates compose this object with the proposition being checked and its
Boolean result. Verification recomputes the two stages from the retained
immutable index, while an explanation query exposes active stamped observations
for exactly the cells consulted by the derivation. Positive, negative, and
tampered certificates are tested. Event-backed states also retain their update
history; snapshot states explicitly report that no history was supplied.

\paragraph{Dependence-aware projection.}
The corruption ablation gives a direct counterexample to treating two nominal
sources as two independent observations. Group-count gates remain unchanged
under the declared correlated copy; raw and source-counting gates do not. The
formal theorem gives the same result for arbitrary finite states satisfying
the group-subset premise.

\paragraph{Model-family answers.}
Robust/permissive relations expose uncertainty that a single matrix suppresses.
Here the envelope covers more true cells and closures, with its width reported.
The point results also show why the family should not be collapsed by fiat:
the robust two-group gate loses badly, while the permissive and exact-consensus
gates improve the descriptive averages.

These contributions use the WM calculus rather than a fixed PLN rule set. PLN
evidence values can instantiate the evidence carrier and supply additional
deduction or revision rules, but none is required for the conservative FCA
projection or the provenance-state laws.

\section{Limitations}
\label{sec:limitations}

\begin{enumerate}
  \item The corruption data are synthetic perturbations of controlled-random
        contexts. They isolate missingness, flips, dependence, and unknowns,
        but do not establish utility in a natural scientific domain.
  \item The analysis is post-run descriptive. The protocol froze models and
        hypotheses, but did not freeze a sampling model, confidence interval,
        or significance test. We therefore report no $p$-values or inferential
        confidence intervals.
  \item Dependence groups are supplied by the modeller. Invariance is exact
        conditional on the grouping; discovering or validating the groups is
        outside the current experiment.
  \item The envelope comparison does not optimize the coverage--width frontier.
        It includes a fixed lower/upper incomplete-context comparator, but not
        a tuned fuzzy, rough, or graded-possibility FCA method. Brier and ECE
        evaluate only interval midpoints as point forecasts, not calibration
        of the set-valued intervals.
  \item Exact lattice enumeration is restricted by a frozen attribute bound.
        Wider cases validate exact queries only. We make no general
        performance claim against specialized FCA engines.
  \item Lean proves the abstract semantics and FCA bridge, while runtime
        equivalence is established by executable differential tests. The C
        implementation is not extracted from Lean and is not end-to-end
        verified.
  \item The query certificates are replayable state-carrying objects, not
        minimal cryptographic or independently checkable proof terms. They can
        expose every consulted cell and its active observations on demand, but
        do not eagerly materialize those explanations. Snapshot states cannot
        reconstruct an update history that was not supplied; event-backed
        states retain one explicitly.
\end{enumerate}

\section{Reproducibility}
\label{sec:reproducibility}

The runtime artifacts live under \code{benchmarks/wm_fca} in the CeTTa
repository \cite{cetta}. The main protocol is
\code{corruption_protocol.json}; the compressed run receipt and descriptive
analysis are under \code{results}. The staged runtime measurements use
\code{ergonomics_protocol.json} and a separate receipt under \code{results}.
The corruption receipt artifact hash is
\WMFCAReceiptHash, and its decompressed content hash is
\WMFCAReceiptContentHash. The incomplete-context extension is identified by
\WMFCAIncompleteProtocolHash; its receipt artifact and content hashes are
\WMFCAIncompleteReceiptHash{} and \WMFCAIncompleteReceiptContentHash.
Exact and query oracles pin source commits, input hashes, dependency versions,
matrix semantics, and oracle agreement.

The Lean source is
\code{Mettapedia/Logic/ConceptOntology/WMFCA.lean} in MeTTapedia
\cite{mettapedia}. Paper tables are emitted by
\code{papers/generate_wm_fca_benchmark_tables.py}. The generator independently
replays both the corruption and incomplete-context analyzers from their pinned
compressed receipts and requires exact JSON equality with the checked-in
analyses. It also checks the protocol hashes, complete and identical case
universes, exact lower-completion/crisp equality, cross-analysis WM means,
oracle schemas and agreement, event-prefix cardinality, and every referenced
field. For the ergonomics table it checks the runner, library, program,
golden-output, and protocol hashes and rederives each reported rate from its
median. A missing or drifted field aborts generation. Thus no numeric table
entry in this paper is silently defaulted or hand-transcribed.

\section{Related work}

FCA originates with Wille \cite{wille1982}; the standard mathematical and
algorithmic treatment is Ganter and Wille \cite{ganterwille1999}. Fuzzy Galois
connections and fuzzy FCA replace Boolean incidence by graded relations
\cite{belohlavek1999}, while interval-valued FCA represents lower and upper
membership degrees \cite{djouadi2009}. The present construction differs in
where uncertainty lives: a provenance state and a family of evidence
extraction models precede the crisp relation. Fuzzy or interval evidence can
nevertheless be used inside the same interface.

Epistemic incomplete contexts distinguish positive, negative, and unknown
incidences and study possible complete contexts obtained by resolving the
unknown cells \cite{chacongomez2026}. Our fixed lower/upper-completion
comparator instantiates that established boundary on exactly the same
corrupted primary layers as the WM models. The WM envelope instead ranges over
explicit evidence profiles and extraction gates, so its bounds need not equal
the extremal completions of one partial matrix.

The WM formulation is aligned with imprecise-probability practice in retaining
a family of admissible precise models instead of prematurely selecting one
\cite{walley1991}. Its exact FCA projection and credal concept carriers extend
the type-theoretic development in \cite{credalconcept}. CeTTa supplies the
native MeTTa execution substrate; Lean~4 supplies the proof-checked semantic
bridge \cite{lean4}.

\section{Conclusion}

World-model FCA can retain observation status, provenance, dependence,
revision, forgetting, and model uncertainty without sacrificing classical
FCA. The crisp theory is recovered exactly as a singleton late projection.
The runtime reconstructs independently computed concepts, implications, and
lattice covers; wider contexts remain safely query-bounded; incremental
updates match fresh recomputation; and the Lean development proves the
corresponding abstract laws.

The frozen corruption study provides a first nontrivial use of the additional
degrees of freedom. Dependence-aware gates eliminate correlated-copy
inflation exactly. Permissive and exact-consensus projections improve several
descriptive recovery metrics, while the robust projection worsens them.
Contextual envelopes improve coverage with an explicit width cost, and their
midpoint point forecasts improve Brier score and ECE in every frozen
missingness and flip-rate stratum. Against the fixed incomplete-context
baseline, WM permissive has better point and structural recovery, whereas the
completion interval has better cell coverage and midpoint calibration at
greater width. These results validate the machinery and isolate its behavior,
but they do not establish interval calibration or natural-domain superiority.
The next scientific step is a preregistered partial-observation dataset in
which provenance and dependence are externally meaningful, together with a
coverage--sharpness comparison against graded uncertain-FCA methods.

\begin{thebibliography}{99}

\bibitem{wille1982}
R.~Wille.
\newblock Restructuring lattice theory: an approach based on hierarchies of concepts.
\newblock In \emph{Ordered Sets}, pp. 445--470. Reidel, 1982.

\bibitem{ganterwille1999}
B.~Ganter and R.~Wille.
\newblock \emph{Formal Concept Analysis: Mathematical Foundations}.
\newblock Springer, 1999.

\bibitem{belohlavek1999}
R.~Belohlavek.
\newblock Fuzzy Galois connections.
\newblock \emph{Mathematical Logic Quarterly}, 45(4):497--504, 1999.

\bibitem{djouadi2009}
Y.~Djouadi and H.~Prade.
\newblock Interval-valued fuzzy formal concept analysis.
\newblock In \emph{ISMIS 2009}, LNCS 5722, pp. 592--601. Springer, 2009.

\bibitem{chacongomez2026}
F.~Chac\'on-G\'omez, M.~E.~Cornejo, D.~Dubois, J.~Medina, and H.~Prade.
\newblock Handling incomplete information in formal concept analysis---a
possibilistic approach.
\newblock \emph{Computational and Applied Mathematics}, 45:171, 2026.
\newblock \url{https://doi.org/10.1007/s40314-025-03440-3}.

\bibitem{walley1991}
P.~Walley.
\newblock \emph{Statistical Reasoning with Imprecise Probabilities}.
\newblock Chapman \& Hall, 1991.

\bibitem{fcar}
M\'alaga FCA Group.
\newblock Concept lattices with \texttt{fcaR}: the Planets context.
\newblock \url{https://malaga-fca-group.github.io/fcaR/articles/concept_lattice.html}.

\bibitem{fcapy}
E.~Dudyrev.
\newblock FCApy benchmarks.
\newblock \url{https://github.com/EgorDudyrev/FCApy_benchmarks}.

\bibitem{bobross}
FiveThirtyEight.
\newblock A statistical analysis of the work of Bob Ross.
\newblock \url{https://github.com/fivethirtyeight/data/tree/master/bob-ross}.

\bibitem{lean4}
L.~de Moura and S.~Ullrich.
\newblock The Lean 4 theorem prover and programming language.
\newblock In \emph{CADE 28}, LNCS 12699, pp. 625--635. Springer, 2021.

\bibitem{credalconcept}
Z.~Goertzel.
\newblock Credal concept formation in dependent type theory:
FCA recovery, Walley conjugacy, and a Mizar benchmark.
\newblock MeTTapedia working paper, 2026.

\bibitem{cetta}
CeTTa contributors.
\newblock CeTTa: a C implementation of MeTTa.
\newblock \url{https://github.com/godelclaw/CeTTa}.

\bibitem{mettapedia}
Mettapedia contributors.
\newblock MeTTapedia formalization library.
\newblock \url{https://github.com/godelclaw/MeTTapedia}.

\end{thebibliography}

\end{document}
